% Fichier engendré par outils/generer-tex.py à partir de 07-nombres-complexes.
% Les démonstrations en français sont transcrites du fichier Lean (skill transcrire-preuve-lean).
\documentclass[11pt,a4paper]{article}

% Compile aussi bien avec pdflatex qu'avec xelatex ou tectonic.
\usepackage{iftex}
\ifPDFTeX
  \usepackage[T1]{fontenc}
  \usepackage[utf8]{inputenc}
\else
  \usepackage{fontspec}
\fi
\usepackage[french]{babel}
\usepackage{amsmath,amssymb,amsthm}
\usepackage[margin=2.5cm]{geometry}
\usepackage[hidelinks]{hyperref}

\theoremstyle{definition}
\newtheorem{definition}{Définition}
\theoremstyle{plain}
\newtheorem{theoreme}{Théorème}
\newtheorem{lemme}[theoreme]{Lemme}

% Renvoi à la déclaration Lean, une ligne après la démonstration, aligné à droite.
\newcommand{\source}[2]{\par\smallskip\noindent\hfill{\footnotesize\href{#1}{\texttt{#2}}}\par\smallskip}

\title{Nombres complexes}
\date{}

\begin{document}
\maketitle
% Les identifiants Lean en machine à écrire ne se coupent pas : on tolère des
% espacements plus lâches plutôt que des lignes qui débordent.
\sloppy

\section{NombresComplexes}

\noindent
Lycée \textemdash{} section \og{} Nombres complexes \fg{}, niveau terminale S.

\medskip\noindent
Dans Mathlib, un nombre complexe \emph{est} un couple de réels : $z.\mathrm{re}$ et
$z.\mathrm{im}$ sont ses parties réelle et imaginaire, et $i$ est le couple $(0, 1)$. La
forme algébrique n'est donc pas une définition mais un théorème, le premier de ce
chapitre. Le conjugué s'écrit $\bar z$, le module $|z|$ \textemdash{} c'est la norme de
$\mathbb{C}$ vu comme plan euclidien \textemdash{} et l'argument $\arg z$.

\medskip\noindent
Un point mérite l'attention : $\arg z$ est un \emph{réel}, choisi dans
$]-\pi\,;\pi]$, alors que l'argument du cours est défini modulo $2\pi$. La formule
$\arg(zz') = \arg z + \arg z'$ est donc fausse telle quelle sur les réels, et l'énoncé
correspondant est écrit dans le groupe des angles, où l'égalité a lieu.

\subsection{Forme algébrique}

\begin{theoreme}
Tout complexe s'écrit sous forme algébrique,
\[ z = \mathrm{Re}(z) + i\,\mathrm{Im}(z) , \]
cette écriture est unique \textemdash{} si $a + ib = 0$ avec $a$ et $b$ réels, alors
$a = b = 0$ \textemdash{} et $i^{2} = -1$.
\end{theoreme}

\begin{proof}
Les deux premiers points sont ceux de la bibliothèque : l'égalité
$z = \mathrm{Re}(z) + i\,\mathrm{Im}(z)$ y est démontrée à partir de la construction de
$\mathbb{C}$ comme couples de réels, et $i^{2} = -1$ de même.

Pour l'unicité, on écrit que deux complexes sont égaux lorsque leurs parties réelles et
leurs parties imaginaires le sont. L'hypothèse $a + ib = 0$ donne donc deux égalités de
réels, $a = 0$ et $b = 0$, qui sont exactement la conclusion.

\medskip\noindent
Cette unicité est ce qui autorise l'\emph{identification} des parties réelles et
imaginaires, geste constant du calcul sur les complexes.
\end{proof}
\source{https://github.com/Commutator-IO/learn-lean/blob/main/cours/02-lycee/07-nombres-complexes/NombresComplexes.lean\#L20}{NombresComplexes.lean\#L20}

\subsection{Conjugué}

\begin{theoreme}
Le conjugué respecte la somme et le produit, et son produit avec $z$ donne le carré du
module :
\[ \overline{z + z'} = \bar z + \bar{z'}, \qquad
   \overline{zz'} = \bar z \,\bar{z'}, \qquad
   z \bar z = |z|^{2} . \]
\end{theoreme}

\begin{proof}
Les deux premières égalités disent que la conjugaison est un morphisme d'anneaux, ce
qu'elle est par construction dans la bibliothèque.

Pour la troisième, la bibliothèque donne $z \bar z = \mathrm{N}(z)$, où
$\mathrm{N}(z) = \mathrm{Re}(z)^{2} + \mathrm{Im}(z)^{2}$ est la \emph{norme
algébrique}, un réel. Il reste à savoir que cette quantité est le carré du module, ce
qu'affirme un autre résultat de la bibliothèque, et à voir l'égalité de réels comme une
égalité de complexes.
\end{proof}
\source{https://github.com/Commutator-IO/learn-lean/blob/main/cours/02-lycee/07-nombres-complexes/NombresComplexes.lean\#L30}{NombresComplexes.lean\#L30}

\begin{theoreme}
Un complexe est réel si et seulement s'il est égal à son conjugué, imaginaire pur si et
seulement s'il est l'opposé de son conjugué :
\[ \mathrm{Im}(z) = 0 \iff \bar z = z , \qquad
   \mathrm{Re}(z) = 0 \iff \bar z = -z . \]
\end{theoreme}

\begin{proof}
Dans les deux cas on écrit l'égalité de complexes comme l'égalité des parties réelles et
des parties imaginaires.

Pour la première : $\bar z = z$ signifie $\mathrm{Re}(z) = \mathrm{Re}(z)$, ce qui ne dit
rien, et $-\mathrm{Im}(z) = \mathrm{Im}(z)$, ce qui équivaut à $\mathrm{Im}(z) = 0$.

Pour la seconde : $\bar z = -z$ signifie $\mathrm{Re}(z) = -\mathrm{Re}(z)$, soit
$\mathrm{Re}(z) = 0$, la condition sur les parties imaginaires étant automatiquement
vérifiée.
\end{proof}
\source{https://github.com/Commutator-IO/learn-lean/blob/main/cours/02-lycee/07-nombres-complexes/NombresComplexes.lean\#L41}{NombresComplexes.lean\#L41}

\subsection{Module}

\begin{theoreme}
\[ |zz'| = |z|\,|z'| , \qquad |z + z'| \le |z| + |z'| . \]
\end{theoreme}

\begin{proof}
Ce sont la multiplicativité de la norme et l'inégalité triangulaire, citées à la
bibliothèque : $\mathbb{C}$ y est un corps normé, et son module en est la norme.
\end{proof}
\source{https://github.com/Commutator-IO/learn-lean/blob/main/cours/02-lycee/07-nombres-complexes/NombresComplexes.lean\#L58}{NombresComplexes.lean\#L58}

\subsection{Argument et forme trigonométrique}

\begin{theoreme}
Forme trigonométrique :
\[ z = |z|\bigl(\cos(\arg z) + i \sin(\arg z)\bigr) . \]
\end{theoreme}

\begin{proof}
C'est la définition de l'argument dans la bibliothèque, où $\arg z$ est construit de sorte
que cette égalité soit vérifiée ; le résultat correspondant y est marqué comme
simplification automatique, si bien que la transcription se réduit à l'invoquer.
\end{proof}
\source{https://github.com/Commutator-IO/learn-lean/blob/main/cours/02-lycee/07-nombres-complexes/NombresComplexes.lean\#L65}{NombresComplexes.lean\#L65}

\begin{theoreme}
Pour $z$ et $z'$ non nuls, l'argument d'un produit est la somme des arguments,
\[ \arg(zz') = \arg z + \arg z' , \]
égalité qui a lieu \emph{modulo} $2\pi$, c'est-à-dire dans le groupe des angles.
\end{theoreme}

\begin{proof}
Résultat de la bibliothèque, cité. L'énoncé mérite qu'on s'y arrête, parce que sa forme
Lean est exactement celle du cours et pas celle qu'on écrirait naïvement.

L'argument de la bibliothèque est un réel choisi dans $]-\pi\,;\pi]$. Avec ce choix,
l'égalité entre réels est fausse : pour $z = z' = e^{3i\pi/4}$, la somme des arguments vaut
$3\pi/2$, qui sort de l'intervalle, tandis que $zz' = e^{3i\pi/2} = -i$ a pour argument
$-\pi/2$. Le théorème est donc énoncé après passage au quotient par $2\pi$, dans le type
des angles réels, où l'égalité vaut sans restriction. C'est la traduction fidèle du
\og{} $[2\pi]$ \fg{} du lycée.
\end{proof}
\source{https://github.com/Commutator-IO/learn-lean/blob/main/cours/02-lycee/07-nombres-complexes/NombresComplexes.lean\#L71}{NombresComplexes.lean\#L71}

\subsection{Forme exponentielle, formules d'Euler et de Moivre}

\begin{theoreme}
\[ e^{i\theta} = \cos\theta + i\sin\theta \qquad (\theta \in \mathbb{R}) . \]
\end{theoreme}

\begin{proof}
Résultat de la bibliothèque, qui l'obtient en regroupant les termes pairs et impairs de la
série de l'exponentielle complexe : les premiers donnent la série du cosinus, les seconds
celle du sinus.
\end{proof}
\source{https://github.com/Commutator-IO/learn-lean/blob/main/cours/02-lycee/07-nombres-complexes/NombresComplexes.lean\#L78}{NombresComplexes.lean\#L78}

\begin{theoreme}
Formules d'Euler :
\[ \cos t = \frac{e^{it} + e^{-it}}{2} , \qquad
   \sin t = \frac{e^{it} - e^{-it}}{2i} . \]
\end{theoreme}

\begin{proof}
On part de l'égalité précédente, écrite pour $t$ puis pour $-t$ :
\[ e^{it} = \cos t + i\sin t , \qquad
   e^{-it} = \cos(-t) + i\sin(-t) = \cos t - i\sin t , \]
la seconde utilisant la parité du cosinus et l'imparité du sinus.

En ajoutant les deux, les termes en $\sin$ s'annulent et il reste $2\cos t$ ; en les
retranchant, ce sont les termes en $\cos$ qui s'annulent et il reste $2i\sin t$. Les deux
formules s'obtiennent en divisant, la seconde par $2i$ qui n'est pas nul.
\end{proof}
\source{https://github.com/Commutator-IO/learn-lean/blob/main/cours/02-lycee/07-nombres-complexes/NombresComplexes.lean\#L83}{NombresComplexes.lean\#L83}

\begin{theoreme}
Formule de Moivre :
\[ (\cos\theta + i\sin\theta)^{n} = \cos(n\theta) + i\sin(n\theta) \qquad (n \in \mathbb{N}) . \]
\end{theoreme}

\begin{proof}
On récrit le membre de gauche sous forme exponentielle : c'est $(e^{i\theta})^{n}$. La
propriété fondamentale de l'exponentielle, $e^{u}$ élevé à la puissance $n$ vaut $e^{nu}$,
donne $e^{i n\theta}$, que l'on récrit sous forme trigonométrique. La formule de Moivre
n'est donc rien d'autre que la règle des exposants, lue dans les deux sens.
\end{proof}
\source{https://github.com/Commutator-IO/learn-lean/blob/main/cours/02-lycee/07-nombres-complexes/NombresComplexes.lean\#L102}{NombresComplexes.lean\#L102}

\subsection{Équation du second degré à discriminant négatif}

\begin{theoreme}
Soient $a \ne 0$, $b$, $c$ réels avec $\Delta = b^{2} - 4ac < 0$. L'équation
$az^{2} + bz + c = 0$ a alors deux racines complexes conjuguées et distinctes,
\[ z = \frac{-b + i\sqrt{4ac - b^{2}}}{2a} \quad\text{et}\quad \bar z . \]
\end{theoreme}

\begin{proof}
Comme $\Delta < 0$, le nombre $4ac - b^{2}$ est strictement positif : on peut poser
$r = \sqrt{4ac - b^{2}}$, qui vérifie $r > 0$ et $r^{2} = 4ac - b^{2}$.

\smallskip\noindent
\emph{Le conjugué.} Conjuguer $z = \dfrac{-b + ir}{2a}$ revient à conjuguer numérateur et
dénominateur ; $a$, $b$, $r$ étant réels, seuls $i$ change de signe :
\[ \bar z = \frac{-b - ir}{2a} . \]

\smallskip\noindent
\emph{$z$ est racine.} On réduit au même dénominateur $4a$ et on développe :
\[ a z^{2} + b z + c
   = \frac{(-b + ir)^{2}}{4a} + \frac{b(-b + ir)}{2a} + c
   = \frac{b^{2} - 2ibr + i^{2}r^{2} - 2b^{2} + 2ibr + 4ac}{4a}
   = \frac{4ac - b^{2} - r^{2}}{4a} , \]
où les termes en $ibr$ se sont éliminés et où $i^{2} = -1$. Le numérateur est nul par
définition de $r$. Le même calcul, avec $-ir$ au lieu de $ir$, montre que $\bar z$ est
racine : les termes en $ibr$ s'éliminent de la même façon.

\smallskip\noindent
\emph{Les deux racines sont distinctes.} Si $z = \bar z$, alors, les deux fractions ayant
le même dénominateur non nul, leurs numérateurs sont égaux, ce qui donne
$4a\,(ir) = 0$. Or $4a \ne 0$ et $ir \ne 0$ puisque $r > 0$ et $i \ne 0$ : contradiction.
\end{proof}
\source{https://github.com/Commutator-IO/learn-lean/blob/main/cours/02-lycee/07-nombres-complexes/NombresComplexes.lean\#L113}{NombresComplexes.lean\#L113}

\subsection{Interprétation géométrique}

\begin{theoreme}
La distance de deux points est le module de la différence de leurs affixes :
\[ AB = |z_B - z_A| . \]
\end{theoreme}

\begin{proof}
La distance de $\mathbb{C}$ est celle du plan euclidien, définie par
$\mathrm{dist}(z, w) = |z - w|$. En échangeant les deux arguments \textemdash{} la
distance est symétrique \textemdash{} on obtient la forme du cours, avec l'affixe
d'arrivée moins celle de départ.
\end{proof}
\source{https://github.com/Commutator-IO/learn-lean/blob/main/cours/02-lycee/07-nombres-complexes/NombresComplexes.lean\#L147}{NombresComplexes.lean\#L147}

\begin{theoreme}
L'angle orienté entre deux vecteurs est l'argument du quotient de leurs affixes : pour
$B \ne A$,
\[ \bigl(\overrightarrow{AB}, \overrightarrow{AC}\bigr)
   = \arg\frac{z_C - z_A}{z_B - z_A} . \]
\end{theoreme}

\begin{proof}
La bibliothèque définit l'angle orienté de deux vecteurs du plan complexe par
$\arg(\bar w z)$, où $w$ et $z$ sont leurs affixes. Il s'agit donc de comparer $\bar w z$
et $z/w$ pour $w = z_B - z_A$ et $z = z_C - z_A$.

Or l'inverse d'un complexe s'écrit $w^{-1} = \bar w / \mathrm{N}(w)$, de sorte que
\[ \frac{z}{w} = \frac{1}{\mathrm{N}(w)} \times \bar w z . \]
Les deux nombres ne diffèrent donc que par le facteur $1/\mathrm{N}(w)$, qui est un réel
strictement positif puisque $w \ne 0$. Comme multiplier par un réel strictement positif ne
change pas l'argument, les deux arguments coïncident.
\end{proof}
\source{https://github.com/Commutator-IO/learn-lean/blob/main/cours/02-lycee/07-nombres-complexes/NombresComplexes.lean\#L151}{NombresComplexes.lean\#L151}

\subsection{Alignement, orthogonalité, cercle}

\begin{theoreme}
Soient $A$, $B$, $C$ d'affixes $a$, $b$, $c$, avec $B \ne A$. Alors :
\[ A, B, C \text{ alignés} \iff \mathrm{Im}\frac{c - a}{b - a} = 0 , \]
\[ (AB) \perp (AC) \iff \mathrm{Re}\frac{c - a}{b - a} = 0 . \]
Autrement dit, le quotient est réel dans le premier cas, imaginaire pur dans le second.
\end{theoreme}

\begin{proof}
\emph{Alignement.} L'alignement est écrit sous la forme : il existe un réel $k$ tel que
$c - a = k(b - a)$.

Si c'est le cas, le quotient vaut $k$ après simplification par $b - a$, qui n'est pas nul ;
sa partie imaginaire est donc nulle.

Réciproquement, si la partie imaginaire du quotient est nulle, ce quotient est égal à sa
propre partie réelle, disons $k$ ; en multipliant par $b - a$ on retrouve $c - a = k(b-a)$,
et $k$ convient.

\smallskip\noindent
\emph{Orthogonalité.} L'orthogonalité de $\overrightarrow{AB}$ et $\overrightarrow{AC}$
s'écrit, en coordonnées, comme la nullité du produit scalaire
\[ (b-a)_{\mathrm{re}}\,(c-a)_{\mathrm{re}} + (b-a)_{\mathrm{im}}\,(c-a)_{\mathrm{im}} = 0 . \]
Or la bibliothèque donne la partie réelle d'un quotient :
\[ \mathrm{Re}\frac{z}{w}
   = \frac{z_{\mathrm{re}} w_{\mathrm{re}} + z_{\mathrm{im}} w_{\mathrm{im}}}{\mathrm{N}(w)} , \]
c'est-à-dire le produit scalaire divisé par $\mathrm{N}(w)$, qui n'est pas nul. Les deux
conditions sont donc équivalentes : un quotient est nul si et seulement si son numérateur
l'est.
\end{proof}
\source{https://github.com/Commutator-IO/learn-lean/blob/main/cours/02-lycee/07-nombres-complexes/NombresComplexes.lean\#L167}{NombresComplexes.lean\#L167}

\begin{theoreme}
Équation cartésienne d'un cercle : pour $r \ge 0$,
\[ |z - \omega| = r \iff (x - x_\omega)^{2} + (y - y_\omega)^{2} = r^{2} , \]
où $x$, $y$ et $x_\omega$, $y_\omega$ sont les parties réelles et imaginaires de $z$ et
$\omega$.
\end{theoreme}

\begin{proof}
Le carré du module se lit sur les coordonnées :
\[ |z - \omega|^{2} = \mathrm{N}(z - \omega)
   = (x - x_\omega)^{2} + (y - y_\omega)^{2} . \]

Si $|z - \omega| = r$, il suffit d'élever au carré. Réciproquement, si la somme des carrés
vaut $r^{2}$, alors $|z-\omega|^{2} = r^{2}$ ; les deux nombres $|z - \omega|$ et $r$ étant
positifs ou nuls, leur égalité s'ensuit \textemdash{} c'est ici que sert l'hypothèse
$r \ge 0$, sans laquelle l'équivalence serait fausse.
\end{proof}
\source{https://github.com/Commutator-IO/learn-lean/blob/main/cours/02-lycee/07-nombres-complexes/NombresComplexes.lean\#L192}{NombresComplexes.lean\#L192}

\subsection{Écriture complexe des transformations}

\begin{theoreme}
Écriture complexe des transformations. La translation $z \mapsto z + b$ conserve les
distances ; la rotation de centre $\omega$ et d'angle $\theta$, $z \mapsto \omega +
e^{i\theta}(z - \omega)$, aussi ; l'homothétie de centre $\omega$ et de rapport $k$,
$z \mapsto \omega + k(z - \omega)$, multiplie les distances par $|k|$.
\end{theoreme}

\begin{proof}
Dans les trois cas on calcule la différence des images.

\emph{Translation :} $(z + b) - (z' + b) = z - z'$, donc les modules sont égaux.

\emph{Rotation :} la différence vaut $e^{i\theta}(z - z')$, dont le module est
$|e^{i\theta}|\,|z - z'|$. Or le module de $e^{i\theta}$ vaut $1$ pour $\theta$ réel
\textemdash{} résultat de la bibliothèque \textemdash{} d'où la conservation des distances.

\emph{Homothétie :} la différence vaut $k(z - z')$, de module $|k|\,|z - z'|$, le module
d'un réel étant sa valeur absolue.

\medskip\noindent
Les trois calculs sont le même : la transformation est affine, et seul le coefficient
multiplicatif compte pour l'effet sur les longueurs.
\end{proof}
\source{https://github.com/Commutator-IO/learn-lean/blob/main/cours/02-lycee/07-nombres-complexes/NombresComplexes.lean\#L210}{NombresComplexes.lean\#L210}

\vfill
\noindent\rule{0.35\linewidth}{0.4pt}\par
\noindent{\footnotesize Leonardo de Moura et Sebastian Ullrich,
\emph{The Lean~4 Theorem Prover and Programming Language},
\emph{Automated Deduction — CADE~28}, Springer, 2021, p.~625--635,
\href{https://doi.org/10.1007/978-3-030-79876-5_37}{doi:10.1007/978-3-030-79876-5\_37}.\par}

\end{document}
